\section{Controller and adjudication method}
\label{sec:method}

\subsection{Policy architectures}

We implement eight policies covering the full spectrum from no verification
to oracle-level information:

\begin{enumerate}
    \item \textbf{Never verify} ($\pi_{\text{never}}$): always select
          \texttt{continue} regardless of belief state.
    \item \textbf{Always verify} ($\pi_{\text{always}}$): select
          \texttt{verify} at every eligible step.
    \item \textbf{Fixed periodic} ($\pi_{\text{fixed}}$): verify at a fixed
          step interval or at the terminal step only, with the interval
          tuned on the development split.
    \item \textbf{Random matched budget} ($\pi_{\text{random}}$): verify
          with probability calibrated to match the expected verification
          budget of the adaptive policy, with the probability fixed before
          evaluation.
    \item \textbf{Threshold on self-confidence} ($\pi_{\text{conf}}$):
          verify when the agent's own reported confidence falls below a
          threshold selected on the development split.
    \item \textbf{Threshold on R1 belief} ($\pi_{\text{thresh}}$): verify
          when $q_t$ exceeds a threshold $\theta$, with $\theta$ selected
          on the development split.
    \item \textbf{Myopic value of information} ($\pi_{\text{voi}}$):
          estimate $\mathbb{E}[\Delta U \mid \texttt{verify}(c)]$ as the
          expected reduction in decision-theoretic loss from observing the
          verifier's outcome, approximated by the KL divergence between
          current and expected posterior belief multiplied by the cost of
          the action the belief would change. Select the action with
          maximum estimated net benefit.
    \item \textbf{Oracle} ($\pi_{\text{oracle}}$): uses hidden fault truth
          to select the theoretically optimal action. Reported only as an
          unattainable upper bound; never included in the primary
          comparison.
\end{enumerate}

The adaptive policy ($\pi_{\text{adaptive}}$) is a learned budget-aware
controller trained from replay data or online interaction. The policy
receives the R1 belief state $b_t$, the resource budget remaining, the
set of available actions with their costs, and the eligibility mask (which
actions are permitted under current constraints). It outputs an action
$a_t$ and a selection propensity $p(a_t \mid b_t)$ used for off-policy
evaluation.

\subsection{Adjudication and evidence independence}

Every controller decision is followed by independent adjudication with five
checks: (1)~did the action execute as intended? (2)~was the immediate fault
removed? (3)~is downstream state consistent? (4)~does the scientific claim
remain valid? (5)~is the verifier independent of the policy and training
reward channel? The adjudication record preserves disagreements, missing
evidence, and ambiguous outcomes rather than collapsing them into binary
accept/reject. A verifier-failure status distinguishes cases where the
verification mechanism itself produced an incorrect accept/reject from cases
where a fault was correctly detected but recovery failed.

\subsection{Budget accounting}

For each trajectory, the cumulative costs $C_v(\tau)$,
$C_{\text{compute}}(\tau)$, $C_{\text{human}}(\tau)$, and $D(\tau)$ are
summed over all decision steps. Budget exhaustion triggers a terminal
record with the budget-exhaustion censoring code. The primary evaluation
compares policies at matched expected cost; the per-trajectory realized
cost is allowed to vary stochastically but reported as a distribution.
